\chapter{State of the Art}
While the previous chapter established the theoretical foundations of the utilized frameworks and hardware platforms, this chapter examines the current state of research and technological developments relevant to this thesis. Implementing a robust person-following behavior on a mobile robot requires a synthesis of various disciplines. Therefore, the following sections review recent advancements in mobile robotic computer vision, analyze existing methodologies for real-time person detection, evaluate algorithmic strategies for robotic target-following, and discuss current ROS-based solutions within the academic and industrial communities.

\section{Computer Vision in Mobile Robotics}
The integration of computer vision (CV) has fundamentally transformed the paradigm of mobile robot navigation. Historically, autonomous mobile robots (AMRs) relied primarily on active distance sensors, such as ultrasonic rangefinders or planar 2D LiDAR scanners, to map environments and avoid collisions. While highly effective for geometric obstacle detection, these sensors lack semantic understanding; they cannot differentiate between a static wall and a dynamic human agent. To overcome these limitations, vision-based perception has emerged as a cornerstone of intelligent robotics, serving as the primary source for comprehensive scene understanding.

Over the past two decades, the application of computer vision in mobile robotics has evolved through distinct technological phases. Early implementations utilized classical feature-extraction techniques — such as optical flow, edge detectors, and scale-invariant feature transforms (SIFT) — to achieve basic localization and visual odometry. However, these traditional algorithms often exhibited poor robustness when subjected to dynamic illumination changes or unstructured environments. The field experienced a significant shift with the commercialization of RGB-Depth (RGB-D) sensors, which allowed robotic architectures to seamlessly fuse high-resolution color frames with spatial point-cloud data in real time. 

Today, the state of the art in mobile robotic vision is heavily dominated by deep learning frameworks and Spatial AI. Modern robotic systems utilize Convolutional Neural Networks (CNNs) and Vision Transformers running directly on embedded edge-computing hardware. This architecture enables mobile platforms to perform low-latency semantic segmentation, object detection, and behavioral prediction. By leveraging vision intelligence, contemporary robots can interpret complex spatial relationships and contextual cues, establishing the technological baseline required for complex human-robot interaction (HRI) tasks, such as autonomous person following.

\section{Person Detection Methods} %{TODO überarbeitn}
Robust human detection is a critical prerequisite for successful human-robot co-existence and interaction. Within the domain of computer vision, algorithms designed to isolate and identify human targets have progressed from rigid hand-crafted feature extractors to highly adaptable deep learning architectures. Understanding this evolution is essential for evaluating the performance tradeoffs regarding accuracy and computational latency in real-time robotic frameworks.

Historically, classical pedestrian detection relied heavily on the combination of Histogram of Oriented Gradients (HOG) and Support Vector Machines (SVM). Introduced by Dalal and Triggs, this methodology utilizes local gradient directions to capture the characteristic silhouette of the human body. While computationally lightweight compared to modern neural networks, hand-crafted features suffer from severe limitations in dynamic environments. They exhibit poor resilience against illumination variances, scale fluctuations, and structural occlusions—such as a robot viewing a person partially hidden behind office furniture. Similarly, early cascade classifiers (e.g., the Viola-Jones framework) proved highly efficient for facial detection but lacked the required versatility to track full-body movements from a moving robotic platform.

The emergence of Deep Convolutional Neural Networks (CNNs) fundamentally shifted the state of the art by replacing hand-crafted algorithms with automated feature learning. Early deep-learning approaches utilized two-stage detectors like Faster R-CNN, which first extract regions of interest and subsequently classify them. Although highly accurate, the significant computational overhead of two-stage pipelines introduces latency, making them impractical for closed-loop robotic control where immediate sensor feedback is mandatory. 

To bridge this gap, modern mobile robotics heavily utilizes one-stage detectors, primarily the You Only Look Once (YOLO) framework. By formulating detection as a single regression problem, YOLO predicts bounding boxes and class probabilities directly from full image frames in a single forward pass. This architectural efficiency enables real-time inference speeds exceeding standard camera frame rates. Furthermore, modern iterations of these deep networks transcend basic bounding-box generation by incorporating multi-task learning, such as simultaneous object tracking and pose estimation. Consequently, modern person detection methods not only confirm the presence of a human target but also provide structural data regarding orientation and posture, establishing a robust foundation for active target-following behaviors.


\section{Person Following Approaches}
Once a human target has been successfully detected and isolated within the perception pipeline, the robotic system must translate this spatial data into continuous, safe, and reliable locomotive commands. Strategies for executing a person-following behavior typically fall into two major paradigms: purely reactive geometric control and comprehensive path-planning-based navigation. 

Reactive approaches represent the traditional baseline for target tracking. These methods generally rely on geometric relationships derived from camera coordinate frames or LiDAR clusters to maintain a fixed relative distance and orientation to the target. Control structures often utilize classical Proportional-Integral-Derivative (PID) loops or velocity-vector mapping to continuously adjust the robot’s linear and angular speeds. For example, the control loop attempts to minimize the error between the current target distance and a predefined safe tracking distance (e.g., 1.5 meters), while simultaneously centering the target within the robot's field of view. While reactive algorithms offer exceptionally low computational latency and immediate feedback responses, they possess a critical flaw: environmental blindness. Because reactive controllers map velocity vectors directly to the target without evaluating global spatial structures, the platform is prone to clipping corners or colliding with intervening obstacles in cluttered indoor environments.

To address the limitations of reactive tracking, modern research heavily focuses on path-planning-based navigation approaches. Instead of treating target tracking as a direct line-of-sight vector problem, these methods integrate the detected human coordinates dynamically into the robot's navigation stack—such as the ROS 2 Navigation (Nav2) ecosystem. The human target is defined as a continuously updating moving goal within a local costmap. Local path planners, utilizing algorithms like the Dynamic Window Approach (DWA) or Timed Elastic Band (TEB), then calculate kinematically feasible trajectories. This architecture ensures that the robot actively follows the user while simultaneously calculating smooth trajectories to detour around obstacles, walls, or other dynamic agents. Although navigation-based frameworks exhibit higher computational overhead and require robust real-time localization, they provide the operational safety and reliability necessary for real-world deployment in human-populated environments.


\section{Existing ROS-based Solutions}
Within the open-source robotics community, numerous packages have been developed to achieve target tracking and human-following behaviors. Analyzing these existing implementations is crucial to identify architectural limitations and highlight the specific engineering contributions of this thesis within the ROS 2 ecosystem.

In the legacy ROS 1 paradigm, foundational packages such as \texttt{turtlebot\_follower} or various laser-based tracking nodes established early benchmarks. These implementations primarily relied on geometric clustering of 2D LiDAR data to isolate human legs, or utilized basic depth thresholding from early RGB-D sensors to track the closest moving entity. However, these classical solutions suffered from significant reliability issues; they lacked semantic awareness, making them highly susceptible to target-loss errors whenever the user walked near static obstacles or when another human crossed the robot's line of sight.

With the industry migration to ROS 2, contemporary solutions have increasingly integrated deep learning inference nodes with the advanced Navigation 2 (\texttt{Nav2}) stack. Community-driven packages frequently leverage standard ROS 2 wrappers for deep learning frameworks, such as OpenCV or specialized AI toolkits, to publish bounding boxes as localized spatial goals. Despite their improved tracking robustness, a critical challenge remains regarding hardware deployment on mobile platforms. Standard ROS 2 vision pipelines often demand substantial computational overhead. Running state-of-the-art neural networks directly on resource-constrained central processing units—such as the TurtleBot 4’s built-in Raspberry Pi 4—frequently results in severe frame-rate drops and operational latency, which compromises the real-time responsiveness of the robot's control loop. To bypass this bottleneck, many existing setups require offloading the AI inference to an external, high-performance workstation or hardware.

%This thesis addresses these architectural limitations by proposing a distributed, network-decentralized ROS 2 pipeline. While advanced spatial AI sensors like the Luxonis OAK-D Pro feature an integrated Visual Processing Unit (VPU) designed to offload deep learning inference directly to the edge, exploring this specialized hardware integration was constrained by the development timeline of this project. Instead, this research demonstrates a robust alternative approach by offloading the entire computational burden—including high-level YOLO inference and spatial localization—to an external workstation equipped with a high-performance graphics card. Linked seamlessly via CycloneDDS, this architecture completely frees the central Single-Board Computer from vision-related tasks. It bridges the gap between high-precision deep learning perception and standalone robotic navigation, establishing a highly responsive, low-latency person-following framework tailored for resource-constrained mobile platforms

